% !TEX program = xelatex
% 电工电子技术 · 第一章第四讲（Chapter One-4）· 精讲讲义（lecture-deepdive 规范）
\documentclass[11pt]{ctexart}
%=====================================================================
% theme-violet.tex —— 讲义精讲 · 紫花主题
% 用法：\documentclass[11pt]{ctexart}
%       \input{assets/theme-violet}
%       \begin{document} ... \end{document}
% 只用 ctexart 自带能力 + xcolor/enumitem/fancyhdr/pifont/amsmath，
% 不依赖 tikz/tcolorbox，保证 MiKTeX 可直接编译。
%=====================================================================

\usepackage[a4paper,left=2.0cm,right=2.0cm,top=2.1cm,bottom=2.3cm,%
            headsep=12pt,footskip=22pt]{geometry}
\usepackage{amsmath,amssymb}
\usepackage{xcolor}
\usepackage{array}
\usepackage{enumitem}
\usepackage{pifont}
\usepackage{fancyhdr}
\usepackage[hidelinks]{hyperref}

%---------------------------------------------------------------- 配色
\definecolor{violet}{RGB}{107,78,143}      % 主色：紫
\definecolor{violetdeep}{RGB}{63,44,92}    % 深紫：标题/强调
\definecolor{petal}{RGB}{183,159,208}      % 花瓣：装饰线、花饰
\definecolor{bandviolet}{RGB}{247,243,251} % 浅紫底：档案/提示框
\definecolor{bandcream}{RGB}{252,249,242}  % 米白底：易错/注意框
\definecolor{rulepurple}{RGB}{217,207,230} % 分隔线
\definecolor{inkgray}{RGB}{92,96,110}      % 次要文字

%---------------------------------------------------------------- 字体
% 有 Noto CJK 用 Noto，没有就退回 ctex 默认（Windows 下为宋体/黑体）
\IfFontExistsTF{Noto Serif CJK SC}{\setCJKmainfont{Noto Serif CJK SC}}{}
\IfFontExistsTF{Noto Sans CJK SC}{\setCJKsansfont{Noto Sans CJK SC}}{}

%---------------------------------------------------------------- 版式
\setlist{nosep,leftmargin=1.9em}
\linespread{1.28}
\setlength{\parskip}{3.5pt}
\setlength{\parindent}{0pt}
\setlength{\tabcolsep}{5pt}
\renewcommand{\arraystretch}{1.28}

%---------------------------------------------------------------- 花饰
% \ding{95} = 五瓣小花，\ding{96} = 星形花；都来自 pifont，必定存在
\newcommand{\orn}{\textcolor{petal}{\ding{95}}}          % 行内小标
\newcommand{\ornstar}{\textcolor{petal}{\ding{96}}}      % 星形花
\newcommand{\flowerrule}{%                                % 居中花饰分隔线
  \par\vspace{7pt}%
  \noindent\makebox[\linewidth]{\color{petal}%
    \rule{0.20\linewidth}{0.6pt}\hspace{8pt}{\small\ding{95}}\hspace{4pt}%
    \raisebox{1.1pt}{\tiny\ding{95}}\hspace{4pt}{\small\ding{95}}\hspace{8pt}%
    \rule{0.20\linewidth}{0.6pt}}%
  \par\vspace{7pt}}

%---------------------------------------------------------------- 标题与节头
\newcommand{\doctitle}[3]{%   #1 主标题  #2 副标题  #3 元信息（可用 \\ 分行）
  \begin{center}
    {\color{petal}\small \orn\enspace #2\enspace \ornstar}\\[7pt]
    {\LARGE\bfseries\color{violetdeep} #1}\\[10pt]
    {\color{petal}\rule{0.55\linewidth}{1pt}}\\[4pt]
    {\small\color{inkgray} #3}
  \end{center}\vspace{2pt}}

\newcommand{\hdr}[2]{%   #1 章节号（如 §0）  #2 标题
  \par\vspace{12pt}%
  \noindent\colorbox{violetdeep}{\color{white}\bfseries\small\;#1\;}\hspace{7pt}%
  {\color{violetdeep}\large\bfseries #2}\par\nobreak
  \vspace{3pt}{\color{petal}\hrule height 1.2pt}\vspace{6pt}}

\newcommand{\subhdr}[1]{%
  \par\vspace{6pt}%
  \noindent{\color{violet}\bfseries #1}\par\vspace{1pt}}

%---------------------------------------------------------------- 提示框
\newcommand{\callout}[1]{% 浅紫底：讲义档案、读法说明
  \par\medskip\noindent
  \fcolorbox{rulepurple}{bandviolet}{\parbox{\dimexpr\linewidth-2\fboxsep-2\fboxrule}{\vspace{2pt}#1\vspace{2pt}}}%
  \par\medskip}

\newcommand{\warnbox}[1]{% 米白底：易错、边界条件
  \par\medskip\noindent
  \fcolorbox{rulepurple}{bandcream}{\parbox{\dimexpr\linewidth-2\fboxsep-2\fboxrule}{\vspace{2pt}#1\vspace{2pt}}}%
  \par\medskip}

\newcommand{\gapitem}[2]{% #1 缺口类型（紫） #2 内容
  \par\vspace{2pt}\noindent
  \hangindent=1.4em\hangafter=1
  {\color{petal}\small\ding{96}}\enspace{\bfseries\color{violetdeep}#1}\enspace #2\par}

% 前提旁注：只放在「需要理解的那一步」后面，给陈述 + 本题代入，不重证。
\newcommand{\aside}[1]{%
  \par\vspace{1pt}\noindent
  \hangindent=1.6em\hangafter=1
  {\small\color{inkgray}%
    \hspace*{0.6em}{\color{petal}$\triangleright$}\hspace{2pt}#1}%
  \par\vspace{2pt}}

%---------------------------------------------------------------- 表格
% 本机 MiKTeX 的 colortbl 与 array 版本不匹配（\rowcolor 会报错），
% 故表格统一用「上粗紫线 + 表头 + 中细线」的三线式，不依赖 colortbl。
\newcommand{\tabtop}{\par\vspace{3pt}{\color{violet}\hrule height 1.2pt}\vspace{1pt}}
\newcommand{\tabmid}{\par\vspace{1pt}{\color{petal}\hrule height 0.5pt}\vspace{1pt}}
\newcommand{\tabbot}{\par\vspace{3pt}{\color{rulepurple}\hrule height 0.5pt}\vspace{8pt}}
\newcommand{\thead}[1]{{\bfseries\color{violetdeep}#1}}

%---------------------------------------------------------------- 页眉页脚
\pagestyle{fancy}
\fancyhf{}
\renewcommand{\headrulewidth}{0.4pt}
\renewcommand{\headrule}{\hbox to\headwidth{\color{rulepurple}\leaders\hrule height \headrulewidth\hfill}}
\fancyhead[L]{\small\color{inkgray}\songti \doctag}
\fancyhead[R]{\small\color{inkgray}\songti \coursetag}
\fancyfoot[C]{\small\color{inkgray}\orn\enspace{\color{violet}\thepage}\enspace\orn}

% 每篇讲义在导言区设这两个值
\newcommand{\doctag}{}
\newcommand{\coursetag}{}
\newcommand{\settags}[2]{\renewcommand{\doctag}{#1}\renewcommand{\coursetag}{#2}}

%---------------------------------------------------------------- 数学记号
\newcommand{\vect}[1]{\boldsymbol{#1}}
\newcommand{\dd}{\,\mathrm{d}}
\newcommand{\ddx}{\frac{\mathrm{d}}{\mathrm{d}x}}
\addtolength{\jot}{2pt}



\settags{C1-4 · 精讲}{电工电子技术}

\begin{document}

\doctitle{1.9 叠加定理与 1.10 等效电源定理}
{独立源为什么要分开算、二端网络凭什么只剩两个数、E1 与戴维南主例的每个数字从哪来}
{依据：课件「Chapter One -4-2026中白上传.pdf」（p1--p45，共 45 页；本讲义页码一律指本册 PDF 物理页，引用他册写成 [C4-3 p12] 的形式）\ \textperiodcentered\ 2026-10-23}

%=============================================================
\hdr{§0}{课件坐标、记号与四处拼接}

\subhdr{0.1 这一册其实是两讲：一块讲「怎么算」，一块讲「能不能只留两个数」}

本册前一半（p1--p21）是 1.9 叠加定理：多个独立源同时作用时，把解拆成「每个源单独作用」的分量之和。
后一半（p22--p45）是 1.10 等效电源定理：一个含源线性二端网络对外看，能不能用一个电压源串一只电阻（戴维南）
或一个电流源并一只电阻（诺顿）顶替。
两半不是并列关系：\textbf{戴维南定理就是叠加定理的一次应用}（§4.1 证明）。课件把两讲排在一起却
一句连接都没有，所以本讲义的骨架是「叠加 $\to$ 置零 $\to$ 边界 $\to$ 等效」（§1--§4），
然后两侧各一个主例题（§5），最后才是其余例子与公式卡。

\par\vspace{2pt}
\noindent\begin{tabular}{@{}p{2.1cm}p{3.0cm}p{6.9cm}p{2.4cm}@{}}
\tabtop
\thead{课件页} & \thead{小节} & \thead{页面内容} & \thead{本讲义}
\\
\tabmid
p1 & 1.9 封面 & 标题页，无正文 & §0
\\
p2 & 1.9.1 定义与用途 & 定义句 + 「Purpose: Simplify the calculation」 + 两幅多电源图 & §1.1
\\
p3 & 1.9.2 Content & 正式表述 + 通用网络图 + 叠加式 & §1.1、§1.4
\\
p4 & 1.9.2 Content & 同一页 + 提问「replaced by their internal resistances」 & §2.1
\\
p5 & 1.9.2 Content & 置零规则：电压源短路、电流源开路 & §2.1
\\
p6--p9 & 1.9.2 Content & 同一张图的置零动画（四帧，无新文字） & §2.2
\\
p10 & 1.9.3 Steps & 三步法（三条英文条目） & §2.3
\\
p11 & 1.9.4 Notice & 四条 Notice（第 (3) 条正文重复第 (2) 条） & §3
\\
p12--p15 & 1.9.5 Example E1 & 三支路：$I'=2\,\mathrm{A}$、$I''=-1\,\mathrm{A}$、$I=1\,\mathrm{A}$ & §5.1--§5.4
\\
p16--p18 & 1.9.5 Example E2 & $I'=3\,\mathrm{A}$、$I''=5\,\mathrm{A}$、$I=8\,\mathrm{A}$ & §7.1
\\
p19--p21 & 1.9.5 Example E3--E5 & 只给电路，\textbf{全册没有答案} & §7.2、§7.3、§7.6
\\
p22 & 1.10.1 Basic idea & 二端网络术语、两种分类、两幅等效图 & §4.0
\\
p23 & 1.10.2 定理陈述 & 戴维南定理一句话表述 + 对照图 & §4.1
\\
p24--p25 & 1.10.2 Terminal & 主例原电路；划掉 $8\,\Omega$、留出端口 a、b & §5.5
\\
p26 & 1.10.2 Key & $U_{eS}=U_{OC}$ & §4.2
\\
p27 & 1.10.2 Key & $R_0=U_{eS}/I_{SC}=U_{OC}/I_{SC}$ + 置零求等效电阻 & §4.2、§4.4
\\
p28 & 1.10.2 Key & 「How to deduce this conclusion?」用本网络算了一遍 & §6
\\
p29--p32 & 1.10.2 Steps & Terminal $\to$ Equivalent resistance $\to$ $R_0=6\,\Omega$ & §5.6
\\
p33--p38 & 1.10.2 Steps & Open-circuit voltage $\to U_{ab}=14\,\mathrm{V}$ $\to$ 等效电路 $\to I=1\,\mathrm{A}$ & §5.7、§5.8
\\
p39--p41 & 1.10.2 习题 E1--E3 & p39 与 p16 同电路；p40 与 p20 同电路；p41 换 12 V/6 $\Omega$ & §7.4、§7.3、§7.7
\\
p42--p45 & 1.10.2 习题 E4 & 同一电路贴四个方法名，四页无一个数 & §7.5
\tabbot
\end{tabular}

\subhdr{0.2 四处拼接（课件自己的缝）}

\gapitem{拼接缺口}{p3 说其他源「replaced by their internal resistances」，p5 却说
「For ideal voltage source, it is like a short circuit」与「For ideal current source, it is like a open circuit」。
两种说法之间缺一句桥：置零置的是\textbf{源值}，内阻留在电路里；理想电压源内阻为 $0$ 所以退化成短路，
理想电流源内阻为 $\infty$ 所以退化成开路。§2.1 把这条桥补上。}

\gapitem{拼接缺口}{p21 讲完叠加、p22 直接跳到等效电源，中间没有一句话说明两者的关系。
其实戴维南等效的两条参数（$U_{OC}$ 与 $R_0$）正是叠加定理在「端口」这一个元件上的应用结果：§4.1 给出证明。
不接这一步，两个定理就只能各自背公式。}

\gapitem{形式改写}{p2 的 \texttt{Definition} 与 p3 的 \texttt{Content} 是同一句话的两次换写
（前者按元件说，后者按支路说），真正的机制——为什么「分别作用再相加」等于「一起作用」——两页都没有。
§1.2 用节点方程两行证完。}

\gapitem{版面}{p5--p9 五页是同一张图的连续动画（只演示置零，无新文字）；p29--p38 是同一电路的十帧连续步骤；
p42--p45 是同一电路只换方法名。§0 的坐标表按去重后的内容计，§10 另备案，不重复计缺口。}

\callout{\textbf{本讲义怎么用}：\\
读第一遍只看 §0 表、§5 两个主例题、§8 公式卡——这三处能让你把题做对；\\
读第二遍再看 §1、§2、§4、§6——这四处解释「为什么可以拆、为什么只剩两个数」，是换一个电路还能不能算对的关键；\\
\textbf{§10 缺口清单}逐条对应课件某页某句，课件哪句话让你卡住，就去那里找。}

%=============================================================
\subhdr{0.3 记号：先统一，后文不再回头}

\par\vspace{1pt}{\small
\noindent\begin{tabular}{@{}p{2.9cm}p{11.6cm}@{}}
\tabtop
\thead{记号} & \thead{约定}
\\
\tabmid
$U_S,\ I_S$ & 独立电压源、独立电流源的\textbf{源值}（课件写 $U_{S1},U_{S2},U_{S3},I_S$）；置零就是令它们为 $0$
\\
$I',\ I''$ & 各独立源单独作用时\textbf{同一条支路}电流的分量，\textbf{与总电流取同一参考方向}
\\
$U_{OC}$ & 端口\textbf{开路电压}（课件写 UOC）；开路即端口电流为零时的端电压
\\
$I_{SC}$ & 端口\textbf{短路电流}（课件写 ISC）；短路即端电压为零时从 $+$ 端流出的电流
\\
$R_0$ & 戴维南（也叫端口）\textbf{等效电阻}（课件写 R0）：内部独立源全部置零、受控源保留后，从端口看进去的电阻
\\
$U_{eS}$ & 等效电压源的\textbf{源值}（课件写 UeS）；它与 $U_{OC}$ 数值相等，但一个是「元件参数」、一个是「测量结果」
\\
$R_L$ & 外接负载电阻（p24 的那只 $8\,\Omega$ 就是 $R_L$）
\\
$i_L$ & 从 $+$ 端\textbf{流出}网络、送进负载的电流；戴维南端口方程按这个方向写
\\
$v,\ i$ & 端口电压与\textbf{流入}网络的端口电流（只出现在 §4.1 的定理陈述里）
\tabbot
\end{tabular}}

\warnbox{\textbf{一个记号陷阱}：$U_{eS}$ 与 $U_{OC}$ 数值相等，但 $U_{OC}$ 是「把负载拿掉、量出来的电压」，
$U_{eS}$ 是「画在等效电路里的电源值」。做题时先说清你要的是哪一个，再代数；
求 $R_0$ 时若把负载 $R_L$ 也算进去，$R_0$ 就变成了「端口总电阻」，与定理里的 $R_0$ 不是一回事（§9 第 (4) 条）。}

\subhdr{0.4 本讲义反复用的一个工具：节点法一行模板}

全册的计算都落在同一句话上：把「节点电压」当未知量，对每个节点写「流出该节点的电流之和为零」。
模板是
\[ \sum_{\text{电阻支路}}\frac{V_{\text{本}}-V_{\text{邻}}}{R}
   +\sum(\text{离开本节点的源电流})=0 . \]
三条纪律：先取一个节点为 $0\,\mathrm{V}$（通常取公共底线）；理想电流源直接把常数项写进「离开/进入」；
理想电压源把两个节点钉成一个条件（差值已知），此时对该两支路围成的「超节点」写一次 KCL。

\aside{本讲义有四节要用它，写法完全相同、只是「邻节点」不同：
§5.4 写一次（E1 上端）、§5.7 写两次（主例左右两块）、§7.2 写三次（E3 的 T、M、L）、§7.6 写四次（E5 四个节点）。
把这张模板抄在草稿纸角上，比记四个公式省事。}

%=============================================================
\hdr{§1}{叠加定理为什么成立：线性性的两行证明}

\subhdr{1.1 课件给了什么（p2、p3、p10）}

课件 p2 的定义是（原文）：
\callout{The voltage across (or current through) an element in a linear electric circuit is the
algebraic sum of the voltages across (or currents through) that element due to each independent
source acting alone.}

p3 又换了一种说法，把「元件」改成「支路」，并写成一条叠加式 $I=I'+I''+I'''+I''''$；
p10 给出三步：逐个源单独作用、重复、代数相加。三页合起来是\textbf{结论 + 步骤}，
缺的正是中间那一步——为什么可以分别算再相加。

\subhdr{1.2 机制：节点方程对源值是线性的}

设网络只含线性电阻（电导 $G_k$）、线性受控源与独立源。对每个节点写 KCL，得到

\[ \mathbf{G}\,\mathbf{v}=\mathbf{J} \]

其中 $\mathbf{G}$ 只由各电导和受控源的控制系数决定，\textbf{与源值无关}；$\mathbf{J}$ 则是各独立源源值的线性组合。
把源值分成「第 $k$ 个源单独作用」的若干组，则

\[ \mathbf{J}=\sum_k \mathbf{J}_k ,\qquad
   \mathbf{v}=\mathbf{G}^{-1}\mathbf{J}=\sum_k \mathbf{G}^{-1}\mathbf{J}_k=\sum_k \mathbf{v}_k . \]

支路电压、支路电流都是节点电压的线性组合，所以它们同样满足「分量相加」。

\begin{itemize}
\item 这里唯一被用到的事实是：\textbf{$\mathbf{G}$ 不随源值变}。这正是「元件线性」的含义——电阻的阻值不随电压电流变；
\item 一旦某个元件非线性（二极管、稳压管、白炽灯丝），它的等效电导成了未知量的函数，$\mathbf{G}=\mathbf{G}(\mathbf{v})$，
      $\mathbf{v}=\mathbf{G}(\mathbf{v})^{-1}\mathbf{J}$ 就不再对 $\mathbf{J}$ 线性，叠加立刻失效。
\end{itemize}

\aside{课件 p2 的 「in a linear electric circuit」 与 p11 的 「all of the components must be linear」
说的就是「$\mathbf{G}$ 与源值无关」这一条；本讲义的推导把它写成了可检验的形式：
凡是元件参数会随工作点变的，都不能用叠加。}

\subhdr{1.3 定理陈述（完整形式）}

\callout{\textbf{叠加定理}。\ 设网络 $N$ 由线性元件（含线性受控源）与独立源组成，共有 $n$ 个独立源。
记 $u^{(k)}$ 为「只保留第 $k$ 个独立源、其余独立源源值置零」时某条支路的电压或电流。
则所有独立源同时作用时，该支路的电压或电流为
\[ u=\sum_{k=1}^{n}u^{(k)} . \]
条件有三条：（1）所有元件线性；（2）受控源属于「线性元件」，不能被置零，必须留在每个分电路里；
（3）各分量必须取\textbf{同一个参考方向}，相加是\textbf{代数}相加（带符号）。}

\subhdr{1.4 最小的数字例子：为什么「线性」这么关键}

一只 $1\,\mathrm{k\Omega}$ 电阻接两个同向电流源，$2\,\mathrm{mA}$ 与 $3\,\mathrm{mA}$。
单独作用时端电压分别是 $2\,\mathrm{V}$ 与 $3\,\mathrm{V}$，同时作用是 $5\,\mathrm{V}=2+3$——这就是线性。
换成一只二极管，$2\,\mathrm{mA}$ 时压降约 $0.70\,\mathrm{V}$、$3\,\mathrm{mA}$ 时约 $0.71\,\mathrm{V}$（对数关系，涨得极慢），
两者相加应是 $1.41\,\mathrm{V}$，而同时作用实测只有 $0.71\,\mathrm{V}$ 左右：因为二极管的「电阻」随电流变，
$\mathbf{G}$ 不再是常数。

这个例子说明：叠加不是「把电源一个个关掉再测」的实验操作，而是\textbf{解的结构}。
分电路里的元件值与原电路不同（源被置零了），但拓扑、参考方向、线性元件都必须一一对应。

\subhdr{1.5 「线性」的两个方向：齐次与可加，用 E1 的数验一次}

「线性」包含两条性质，课件只用了其中一条。齐次性：源值同时乘 $\lambda$，所有响应也乘 $\lambda$。
可加性：两个源各自单独作用的结果相加，等于它们同时作用的结果——这就是叠加定理。
合起来才叫线性；只有可加、没有齐次的说法是不成立的（两者其实是同一个线性映射的两面）。

用 E1 验一次。$4\,\mathrm{A}$ 与 $20\,\mathrm{V}$ 时 $I=1\,\mathrm{A}$（§5.4）。把两个源值同时翻倍
（$8\,\mathrm{A}$、$40\,\mathrm{V}$），按齐次性应有 $I=2\,\mathrm{A}$。直接解：上端 KCL 为
$8=\frac{V_T}{10}+\frac{V_T+40}{10}$，得 $V_T=20\,\mathrm{V}$，故 $I=20/10=2\,\mathrm{A}$，与齐次性预期一致。
源值减半（$2\,\mathrm{A}$、$10\,\mathrm{V}$）则给 $I=0.5\,\mathrm{A}$。

这三个数连起来看比单个结果更能说明「线性」是什么意思：\textbf{比例关系对电压、电流成立，对功率不成立}
（功率按平方走，见 §3.3）。做题时用齐次性验量级，比重新解一遍快。

%=============================================================
\hdr{§2}{「置零」到底置掉了什么：规则与三条边界}

\subhdr{2.1 从「内阻」到「短路/开路」：p3 与 p5 怎么对上}

课件 p4 的提问是（原文）：
\callout{What do you mean by saying all the other independent sources are replaced by their
internal resistances? How to deal with the other sources when they are not working?}

p5 给出的答案是两条规则，但没有说明它与 p3 的「internal resistances」是什么关系。桥在这里：

\par\vspace{2pt}
\noindent\begin{tabular}{@{}p{2.6cm}p{4.2cm}p{7.6cm}@{}}
\tabtop
\thead{源} & \thead{真实模型} & \thead{令源值为零之后}
\\
\tabmid
电压源 & 理想源 $U_S$ 串内阻 $R_S$ & 只剩 $R_S$。理想电压源 $R_S=0$，两端被短接 $\Rightarrow$ \textbf{短路}
\\
电流源 & 理想源 $I_S$ 并内阻 $R_S$ & 只剩 $R_S$。理想电流源 $R_S=\infty$，支路断开 $\Rightarrow$ \textbf{开路}
\tabbot
\end{tabular}

所以 p3 的「replaced by their internal resistances」是\textbf{真实源}的说法，p5 的「short / open」是
\textbf{理想源}的说法，两者是同一件事在两个模型下的样子。本册题目里的源一律按理想源处理：
置零 $=$ 电压源短路、电流源开路。

\aside{代入本册的题：E1（p12）只有「$4\,\mathrm{A}$ 理想电流源」与「$20\,\mathrm{V}$ 理想电压源」两个源。
只留 $20\,\mathrm{V}$ 时，$4\,\mathrm{A}$ 那条支路整条断开（不是接 $0\,\mathrm{A}$ 而是开路），
左支路的 $10\,\Omega$ 随之退出电路——所以分电路里那只 $10\,\Omega$ 不起作用。}

\subhdr{2.2 分电路不是「关电源再测」}

p5--p9 用五页动画演示「只让一个源工作」，容易读成「先把别的源关掉，量一下，再换一个」。
准确的说法是：\textbf{把同一个解写成两个子问题的和}。两个子电路是数学上的辅助电路，
它们的源值不同，但节点编号、元件位置、参考方向必须与原电路一一对应，否则分量加不回去。

\warnbox{\textbf{一个常被忽略的边界}：受控源（如三极管的 $\beta i_b$、运放的 $A u_d$）是\textbf{因变量}，
不是独立源。它们\textbf{不能置零}，而且每个分电路里都要保留——因为控制量在每个分电路里都不同。
课件 Notice 第 (1) 条只说了「元件必须线性」，没有点明这一条，而它是叠加法在电子电路里最容易出错的地方。}

\subhdr{2.3 三步法的执行细节（p10）}

p10 的三条（原文）是：
\callout{(1) Turn off all the other independent sources except the one in operation. Find the outputs
(voltage or current) due to that active source.\\
(2) Repeat step (1) for each of the other independent source.\\
(3) Find the total contribution by adding algebraically all the contributions due to the independent sources.}

执行时补三个细节，课件没写但每题都要用：

\begin{enumerate}
\item \textbf{先定参考方向，再拆}。全部分量沿用原电路待求量的箭头方向；分量算出负数就带负号加，
      不要临时改箭头（§3.2 用 E1 的 $I''=-1\,\mathrm{A}$ 演示）。
\item \textbf{一次只留一个源，不等于一次只能留一个}。p11 第 (4) 条 「Each group can have more than one sources」
      允许把若干源合成一组处理（比如两个同向电流源先并联），只要分组不重不漏。
\item \textbf{每一步都要能回到原电路自检}。最省事的自检是把总结果代回原电路直接解一遍（§5.4、§5.8 都这么做了）。
\end{enumerate}

分组的一个小例子：两只 $2\,\mathrm{A}$ 与 $3\,\mathrm{A}$ 的电流源并接在同一对节点上、共同向 $1\,\Omega$ 供电。
分开算得 $2\,\mathrm{V}+3\,\mathrm{V}=5\,\mathrm{V}$；先把两只源合并成 $5\,\mathrm{A}$ 再算，得 $5\,\mathrm{V}$——
同一个数，但少了一个分电路。分组只对\textbf{同类、同向}的源方便，方向不一致时必须带符号合并。

%=============================================================
\subhdr{2.4 受控源不能置零：一个 $6\,\mathrm{V}$ 变 $3\,\mathrm{V}$ 的反例}

上面第 (2) 条边界值得单独算一遍。取一个只有三条支路的电路：节点 A 经 $1\,\Omega$ 接地；
一只 $2\,\mathrm{A}$ 独立电流源、一只 $1\,\mathrm{A}$ 独立电流源、一只压控电流源（$g_m=0.5\,\mathrm{S}$，
控制量取 $v_A$ 本身），三支都由地注入 A。节点法一行：

\[ \frac{v_A}{1}=2+1+0.5\,v_A\ \Longrightarrow\ v_A=6\,\mathrm{V} . \]

正确的叠加：只留 $2\,\mathrm{A}$ 时由 $\frac{v_A}{1}=2+0.5v_A$ 得 $v_A=4\,\mathrm{V}$；
只留 $1\,\mathrm{A}$ 时同法得 $v_A=2\,\mathrm{V}$；相加 $4+2=6\,\mathrm{V}$，与上面一致。
注意每个分电路里那条受控源都在，而且它贡献的量在每个分电路里不同
（$v_A=4\,\mathrm{V}$ 时是 $2\,\mathrm{A}$，$v_A=2\,\mathrm{V}$ 时是 $1\,\mathrm{A}$）。

错误做法：把受控源也当独立源置零，就只剩 $\frac{v_A}{1}=2+1$，$v_A=3\,\mathrm{V}$——正好少一半。
这个 $2$ 倍的来历就是 $\frac{1}{1-0.5}=2$：受控源把结果放大了 $2$ 倍，删掉它就等于把放大能力删掉。

\aside{本册的两个主例（E1 与戴维南主例）都不含受控源，所以这条边界在课件里看不出来；
但本课程第四章的微变等效电路（[C4-3 p12--p18]）里三极管就是受控源，那里一旦置零受控源，
$A_u$ 这种量会整个消失。}

%=============================================================
\hdr{§3}{p11 的四条 Notice：两条能用、一条是空的、一条要补}

\subhdr{3.1 四条的对照}

\par\vspace{2pt}
\noindent\begin{tabular}{@{}p{1.1cm}p{5.2cm}p{8.2cm}@{}}
\tabtop
\thead{条目} & \thead{原文要点} & \thead{本讲义的处理}
\\
\tabmid
(1) & 「all of the components must be linear」 & 见 §1.2；非线性元件（二极管）上不能用叠加
\\
(2) & 「It is not applicable to power」 & 结论对，但课件没给证据；§3.3 补出交叉项与数字
\\
(3) & 标题 「Pay attention to the direction」，正文却整段重复第 (2) 条 & \textbf{这一条是空的}；§3.2 补出方向规则
\\
(4) & 「Each group can have more than one sources」 & 允许分组，见 §2.3 第 (2) 点
\tabbot
\end{tabular}

第 (3) 条值得单独说一句：它的标题承诺讲「方向」，正文却把第 (2) 条关于功率的话一字不差重复了一遍
（p11 上确实如此），所以「注意方向」到底注意什么，课件全程没有答案。下面补上。

\subhdr{3.2 方向规则：分量带符号}

\callout{\textbf{方向规则}。所有分量都在\textbf{原电路的参考方向}下计算；
若某个分量算出来是负数，说明该源单独作用时实际电流方向与参考方向相反，相加时\textbf{保留负号}。
不要因为「看起来是反的」就提前把箭头翻过来——翻箭头等于把符号藏进图里，最后一定错。}

用 E1（p14--p15）的数字看：$I'=2\,\mathrm{A}$、$I''=-1\,\mathrm{A}$，所以 $I=2+(-1)=1\,\mathrm{A}$。
若把 $I''$ 的箭头提前翻成「$1\,\mathrm{A}$ 且方向相反」，再加就变成 $2+1=3\,\mathrm{A}$，错。
反号的来源在 §5.3 讲清楚：$20\,\mathrm{V}$ 源的 $+$ 端在下，它把电流\textbf{顶}回上母线，与 $I$ 的参考方向相反。

\subhdr{3.3 功率为什么不能叠加：交叉项}

p11 第 (2) 条说功率不能用叠加，理由是「power is not a linear quantity」。这句是对的，但读者往往
「知道不对、不知道差多少」。把 $I=I'+I''$ 代进 $P=I^2R$：

\[ P=(I'+I'')^2R=\underbrace{I'^2R}_{P'}+\underbrace{I''^2R}_{P''}+\,2I'I''R . \]

最后那一项就是交叉项，也是「不能叠加」的全部差额。用 E1 的中间支路（$R=10\,\Omega$）算一遍：

\par\vspace{2pt}
\noindent\begin{tabular}{@{}p{4.6cm}p{3.2cm}p{7.0cm}@{}}
\tabtop
\thead{量} & \thead{数值} & \thead{说明}
\\
\tabmid
真值 $P=I^2R$ & $(1\,\mathrm{A})^2\times10\,\Omega=10\,\mathrm{W}$ & 两源同时作用
\\
分量 $P'=I'^2R$ & $(2\,\mathrm{A})^2\times10\,\Omega=40\,\mathrm{W}$ & 只留 $4\,\mathrm{A}$ 源
\\
分量 $P''=I''^2R$ & $(1\,\mathrm{A})^2\times10\,\Omega=10\,\mathrm{W}$ & 只留 $20\,\mathrm{V}$ 源
\\
交叉项 $2I'I''R$ & $2\times2\times(-1)\times10=-40\,\mathrm{W}$ & 负号，因为两分量方向相反
\\
$P'+P''$ & $50\,\mathrm{W}$ & 与真值 $10\,\mathrm{W}$ 差 $40\,\mathrm{W}$，差额正是交叉项
\tabbot
\end{tabular}

所以「功率不能叠加」不是玄学：$40+10-40=10$，交叉项必须补上，而交叉项里有两个源的乘积，
不是任何单个源的贡献。题目若问功率，先求总电流或总电压，再平方——这是唯一不会错的路线。

%=============================================================
\hdr{§4}{两讲其实是同一件事：戴维南定理是叠加定理的一次应用}

\subhdr{4.0 课件 p22 想要什么}

p22 的版面把三件事一次摆出来：术语（「Two terminal network」与有源、无源二端网络）、
分类（「Thévenin's Theorem」「Norton's Theorem」）、用途
（「Simplify an active linear two-terminal network」）。它承诺「把一个有源线性二端网络化简」，
但没有说清「凭什么能化简成两个元件」。p23 直接给出定理陈述：

\callout{Thévenin's theorem states that a linear two-terminal network can be replaced by an
equivalent circuit consisting of a voltage source $U_{OC}$ in series with a resistor $R_0$.}

这句话只说「可以」，没说「为什么可以」。下面把它还原成原形式并证明。

\subhdr{4.1 定理（原形式）与证明}

\callout{\textbf{戴维南定理（原形式）}。\ 设 $N$ 是由线性电阻、线性受控源与独立源组成的二端网络，
端口电压 $v$ 与端口电流 $i$ 取关联参考方向（$i$ 自 $+$ 端\textbf{流入}网络）。则存在两个常数
$U_{OC}$ 与 $R_0$，使对任意 $i$ 都有
\[ v=U_{OC}+R_0\,i . \]
其中 $U_{OC}$ 是端口开路（$i=0$）时的端电压，$R_0$ 是内部独立源全部置零、受控源保留后从端口看进去的等效电阻。\\
若记作「送进负载的电流」$i_L=-i$，则常用形式为
\[ v=U_{OC}-R_0\,i_L . \]
}

证明部分是本讲义补的 [教材补充]：课件 p23 只给陈述，p26--p28 只给两个公式，全册没有一处证明。
补它的理由很实在——不证，「为什么只要两个数」就只能靠背；证明了，这句话就是 §1.2 的直接推论。

\textbf{证明（用叠加定理）}。把端口电流 $i$ 看成「外接一个电流源」的作用。
由 §1.2，端口电压对「内部独立源」与「外接电流源」是线性的，故可拆成两项：
$v=v_{\mathrm{in}}+v_{\mathrm{ex}}$，其中 $v_{\mathrm{in}}$ 是只留内部源时的端口电压，
$v_{\mathrm{ex}}$ 是只留外接源时的端口电压。

第一项里 $i=0$（外接源被置零 $=$ 开路），正是端口开路电压 $U_{OC}$。
第二项里内部独立源全部置零，网络退化成纯电阻网络，$R_0$ 就是它从端口看进去的等效电阻；
此时电流 $i$ 自 $+$ 端流入、从 $-$ 端流出，端电压就是 $v=R_0\,i$。
两项相加即得 $v=U_{OC}+R_0\,i$。

\aside{这一步的旁注只有一句：\textbf{「线性 $\Rightarrow$ 端电压是端口电流的一次函数 $\Rightarrow$ 两个数就够了」}。
对本题的代入选法：主例（p24）里 $U_{OC}=14\,\mathrm{V}$、$R_0=6\,\Omega$，
接 $8\,\Omega$ 负载时 $i_L=1\,\mathrm{A}$，于是 $v=14-6\times1=8\,\mathrm{V}$——与直接量 $8\,\Omega$ 上的电压一致（§5.8 自检）。}

\subhdr{4.2 课件的两个公式是推论，不是定理}

p26 写 「$U_{eS}=U_{OC}$」，p27 写 $R_0=U_{eS}/I_{SC}=U_{OC}/I_{SC}$，并配一句
「Shut down all the sources inside and calculate the equivalent resistance」。
这两条都能从 §4.1 的直线直接读出：

\begin{itemize}
\item 令 $i_L=0$（端口开路）：$v=U_{OC}$。这就把「等效电源的源值」钉死为开路电压（p26）；
\item 令 $v=0$（端口短路）：$i_L=U_{OC}/R_0$。这就是短路电流 $I_{SC}$，反解得 $R_0=U_{OC}/I_{SC}$（p27）；
\item p27 那句「置零内部源、算等效电阻」是同一件事的第二种求法：$R_0$ 本来就是直线斜率，斜率与源值无关，
      所以可以先把源置零再量电阻。
\end{itemize}

顺带说清 $U_{eS}$ 与 $U_{OC}$ 的关系：数值相等，含义不同（§0.3 的记号陷阱）。

\subhdr{4.3 诺顿：同一条直线的另一种写法}

把 $v=U_{OC}-R_0i_L$ 两边除以 $R_0$ 并移项：

\[ i_L=\frac{U_{OC}}{R_0}-\frac{v}{R_0}=I_{SC}-\frac{v}{R_0} . \]

这就是「一个电流源 $I_{SC}$ 并一只 $R_0$」的端口方程（p22 右下图）。
所以戴维南与诺顿不是两个定理，而是\textbf{同一条直线的两种参数化}：一个用截距（$U_{OC}$）与斜率（$-R_0$），
一个用纵截距（$I_{SC}$）与斜率（$-1/R_0$）。做题时任选一种，只要求出的两个数对得上：

\[ U_{OC}=I_{SC}\,R_0 . \]

\subhdr{4.4 $R_0$ 的两种求法与等价条件}

\par\vspace{2pt}
\noindent\begin{tabular}{@{}p{3.2cm}p{5.6cm}p{5.6cm}@{}}
\tabtop
\thead{求法} & \thead{做法} & \thead{什么时候用}
\\
\tabmid
置零求电阻 & 内部独立源置零（电压源短路、电流源开路），受控源保留，从端口看进去算等效电阻 & 网络里电阻都能串并联时最快；含受控源时用它
\\
比例求电阻 & 先求 $U_{OC}$，再求 $I_{SC}$，$R_0=U_{OC}/I_{SC}$ & 网络里有电桥、不好串并联时；但 $R_0=0$ 或 $R_0=\infty$ 时失效（§6.3）
\tabbot
\end{tabular}

两种方法必须给出同一个 $R_0$，这本身就是一道最省时间的自检题。用主例（p24）把两条路并排走一遍：

\par\vspace{2pt}
\noindent\begin{tabular}{@{}p{2.4cm}p{5.4cm}p{5.4cm}@{}}
\tabtop
\thead{步骤} & \thead{置零法（§5.6）} & \thead{比例法（§6.5）}
\\
\tabmid
先求什么 & 不求电压电流，直接把源置零 & 先求 $U_{OC}=14\,\mathrm{V}$（p36）
\\
再求什么 & 从 a、b 看进去算串并联 & 再求 $I_{SC}=7/3\,\mathrm{A}$（短路端口）
\\
$R_0$ & $(2\parallel2)+5=6\,\Omega$ & $14/(7/3)=6\,\Omega$
\\
适用 & 电阻能串并联时最快 & 有电桥、不好串并联时；$R_0=0$ 或 $\infty$ 时失效（§6.3）
\tabbot
\end{tabular}

两条路的中间量完全不同、终点必须相同。这就是「$R_0$ 是端口属性，与你怎么算无关」的实证。

%=============================================================
\subhdr{4.5 边界：非线性网络为什么不能只用两个数}

p23 的前提是「a linear two-terminal network」。换成非线性元件，这条前提就断了，而且这不是抽象的担心：
把端口上的负载电阻换成一个二极管再串一只电阻，端口特性就是一条折线（正向导通段很陡、反向截止段很平），
不是直线。此时用任意两点连一条线去「等效」，第三点必然对不上：
例如用「开路点」与「短路点」定出 $U_{OC}$ 与 $I_{SC}$，代回一个中间工作点，算出来的电压与实测差得很远。
所以「两个数」的办法不是近似技巧，而是线性这一前提的直接推论；前提不在，结论就不在。

同理，若网络里某只电阻的阻值随温度或电压变（白炽灯丝、热敏电阻），端口特性也会弯，
叠加与戴维南会同时失效——它们其实依赖的是同一个条件（§1.2 的 $\mathbf{G}$ 与源值无关）。

%=============================================================
\hdr{§5}{两个主例题全程：E1（p12--p15）与戴维南主例（p24--p38）}

\subhdr{5.1 E1 的电路与分电路（p12--p13）}

p12 的电路是三支路并联在上下两条母线之间，课件只给了数值与箭头，没有给支路清单。
按页把这些元件列成表，后面每个数都从这张表出：

\par\vspace{2pt}
\noindent\begin{tabular}{@{}p{1.4cm}p{8.4cm}p{2.4cm}p{2.4cm}@{}}
\tabtop
\thead{支路} & \thead{自上而下的元件} & \thead{上端} & \thead{下端}
\\
\tabmid
左 & $10\,\Omega$ 串 $4\,\mathrm{A}$ 电流源（源在下，箭头朝上） & T & B
\\
中 & $10\,\Omega$；$I$ 的参考方向取 T$\to$B & T & B
\\
右 & $10\,\Omega$ 串 $20\,\mathrm{V}$ 电压源（源在下，$-$ 在上、$+$ 在下） & T & B
\tabbot
\end{tabular}

把下端 B 取作 $0\,\mathrm{V}$，上端记 T。p13 把原电路拆成两个分电路：右上角只留 $4\,\mathrm{A}$（$20\,\mathrm{V}$ 短路），
右下角只留 $20\,\mathrm{V}$（$4\,\mathrm{A}$ 开路）。

\aside{这里藏着一个读图自检：p12 文字层一共标了\textbf{三只} $10\,\Omega$（每支路一只）。
若把中间支路误读成两只 $10\,\Omega$ 串联，分电路会算出 $4/3\,\mathrm{A}$ 与 $-2/3\,\mathrm{A}$，
与课件给的 $2\,\mathrm{A}$、$-1\,\mathrm{A}$ 都对不上。只有「三只 $10\,\Omega$」这一读法能同时复现课件的两个分量与总数，
可以用它反过来核对电路图。}

\subhdr{5.2 只留 $4\,\mathrm{A}$：$I'=2\,\mathrm{A}$}

$20\,\mathrm{V}$ 源置零即短路，于是右支路只剩 $10\,\Omega$。上端 T 的 KCL 只有两条路可走（中支路与右支路），
二者都是 $10\,\Omega$，所以 $4\,\mathrm{A}$ 平分：

\[ I'=4\times\frac{10}{10+10}=2\,\mathrm{A} , \]

方向由 T 到 B，与参考方向相同。顺带得此时 $V_T=2\,\mathrm{A}\times10\,\Omega=20\,\mathrm{V}$
（与 §5.4 的直接解一致，可作自检）。

\subhdr{5.3 只留 $20\,\mathrm{V}$：$I''=-1\,\mathrm{A}$ 与负号的来历}

$4\,\mathrm{A}$ 源置零即开路，左支路整条断开（那只 $10\,\Omega$ 退出电路）。剩下的是一个单回路：
中支路 $10\,\Omega$ 与右支路（$10\,\Omega$ 串 $20\,\mathrm{V}$）。回路总电阻 $20\,\Omega$、电动势 $20\,\mathrm{V}$，
所以回路上电流大小是 $1\,\mathrm{A}$。方向是关键：右支路电源的 $+$ 端\textbf{在下}，
源内部电流由 $-$ 流向 $+$，即它把电流\textbf{顶向}底母线，再从中支路\textbf{由下往上}回到上母线。

所以中支路的实际电流是从 B 流向 T，与参考方向 T$\to$B 相反：

\[ I''=-\frac{20}{10+10}=-1\,\mathrm{A} . \]

\aside{负号不是算错：它记录的是「该源单独作用时电流方向与参考方向相反」。
若把 $I''$ 的箭头提前翻过来写成 $+1\,\mathrm{A}$，加的时候就会变成 $3\,\mathrm{A}$（§3.2 的反例）。}

\subhdr{5.4 相加与自检：$I=1\,\mathrm{A}$}

\[ I=I'+I''=2+(-1)=1\,\mathrm{A} . \]

p15 的黄色高亮块给的正是这个数。自检用原电路直接解：上端 T 的 KCL（$4\,\mathrm{A}$ 流入 T）为

\[ 4=\frac{V_T}{10}+\frac{V_T+20}{10}\ \Longrightarrow\ V_T=10\,\mathrm{V},\qquad
   I=\frac{V_T}{10}=1\,\mathrm{A} . \]

两边一致，说明「三只 $10\,\Omega$」的读法与课件的三个数全部对得上。

\subhdr{5.5 戴维南主例：端口怎么划（p24--p25、p30）}

p24 的电路有五条支路，$8\,\Omega$ 上标着待求电流 $I$。先把元件列清（B 为底部公共线）：

\par\vspace{2pt}
\noindent\begin{tabular}{@{}p{2.5cm}p{7.2cm}p{2.6cm}p{2.6cm}@{}}
\tabtop
\thead{支路} & \thead{元件} & \thead{上端} & \thead{下端}
\\
\tabmid
左（$8\,\mathrm{V}$） & $8\,\mathrm{V}$ 电压源（$+$ 在上）串 $2\,\Omega$ & L & B
\\
左（竖 $2\,\Omega$） & $2\,\Omega$ & L & B
\\
待求支路 & $8\,\Omega$；$I$ 的参考方向取 a$\to$b & L（$=$a） & R（$=$b）
\\
右（$3\,\mathrm{A}$） & $3\,\mathrm{A}$ 电流源（箭头向下，电流 b$\to$B） & R（$=$b） & B
\\
右（$5\,\mathrm{V}$） & $5\,\mathrm{V}$ 电压源（$+$ 在上）串 $5\,\Omega$ & R（$=$b） & B
\tabbot
\end{tabular}

p25/p30 做的事只有一件：把 $8\,\Omega$ 拿掉，在断口两端标 a、b。这样做合法的理由（p34 的浅蓝方框指的就是它）：
$8\,\Omega$ 是 L 与 R 之间（除公共底线以外）\textbf{唯一}的连接元件，拿掉之后网络分成左右两块，
\textbf{两块只经底部公共线相通}。公共线是参考节点，所以空载时两块之间没有电流交换，
各自可以当「单节点电路」单独解——这是后面所有计算的支点。

\subhdr{5.6 $R_0=6\,\Omega$（p31--p32）}

置零内部源：$8\,\mathrm{V}$ 短路、$5\,\mathrm{V}$ 短路、$3\,\mathrm{A}$ 开路。从 a、b 看进去：

\begin{itemize}
\item 左块：$8\,\mathrm{V}$ 支路的 $2\,\Omega$ 与并联的 $2\,\Omega$ 变成 $2\parallel2=1\,\Omega$（a 到公共线）；
\item 右块：$5\,\mathrm{V}$ 短路后只剩 $5\,\Omega$（公共线到 b）；
\item 两者\textbf{串联}：从 a 到 b 的唯一通路是 a $\to$ 左块 $\to$ 底部公共线 $\to$ $5\,\Omega$ $\to$ b。
\end{itemize}

\[ R_0=(2\parallel2)+5=1+5=6\,\Omega , \]

与 p32 的 $R_0=6\,\Omega$ 一致。

\aside{代入本题：这里最容易错成 $(2\parallel2)\parallel5=0.83\,\Omega$。
判据是「换个端口还能不能走通」——a、b 分别在左块与右块的上端，要从 a 到 b 必须串穿两块，
所以是串联。串联的依据是路径，不是画法。}

\subhdr{5.7 $U_{OC}=14\,\mathrm{V}$（p33--p36），以及 p36 里 5、15、4 三个数}

空载时两块解耦（§5.5），各自一次 KCL。

\textbf{（1）左块：}a 端只有两条路（$8\,\mathrm{V}$ 串 $2\,\Omega$、另一只 $2\,\Omega$），电流之和为零：

\[ \frac{V_a-8}{2}+\frac{V_a}{2}=0\ \Longrightarrow\ V_a=4\,\mathrm{V} . \]

等价地看成一次分压：$8\,\mathrm{V}$ 落在 $2\,\Omega$ 与（$2\,\Omega$）之间，$V_a=8\times2/(2+2)=4\,\mathrm{V}$。

\textbf{（2）右块：}$3\,\mathrm{A}$ 由 b 灌入电流源、向下流到底部；b 只剩另一条路（$5\,\mathrm{V}$ 串 $5\,\Omega$），
所以 $5\,\Omega$ 里也是 $3\,\mathrm{A}$，但方向\textbf{向上}（b 处进去的电流只能从这条支路补）。
沿「从底母线 $0\,\mathrm{V}$ 往上」走：电流向上，与行进方向相同，经过 $5\,\Omega$ 电位\textbf{降} $3\times5=15\,\mathrm{V}$，
于是 $5\,\Omega$ 上端为 $-15\,\mathrm{V}$；再经过 $5\,\mathrm{V}$ 源（由 $-$ 到 $+$）电位\textbf{升} $5\,\mathrm{V}$，故

\[ V_b=0-15+5=-10\,\mathrm{V} . \]

（用方程验证同一件事：设 $5\,\Omega$ 上端为 $x$，向上电流 $3\,\mathrm{A}$ 意味着 $x-0=-3\times5=-15\,\mathrm{V}$，
再由 $x$ 过源升 $5\,\mathrm{V}$ 得 $V_b=-10\,\mathrm{V}$。两条路一致。）

\textbf{（3）开路电压：}

\[ U_{ab}=V_a-V_b=4-(-10)=14\,\mathrm{V} , \]

与 p36 的 $U_{ab}=14\,\mathrm{V}$ 一致。

\textbf{（4）p36 那个式子怎么读。}课件写的是 $U_{ab}+5-15-4=0$，即 $U_{ab}=4+15-5$。
三个数各有出处，课件一个都没标：

\par\vspace{2pt}
\noindent\begin{tabular}{@{}p{2.2cm}p{3.4cm}p{8.6cm}@{}}
\tabtop
\thead{p36 中的数} & \thead{来源} & \thead{说明}
\\
\tabmid
$4$ & $V_a$ & $8\,\mathrm{V}$ 在 $2\,\Omega\parallel2\,\Omega$ 上的分压，即左块的单节点解
\\
$15$ & $3\,\mathrm{A}\times5\,\Omega$ & $3\,\mathrm{A}$ 全部流过 $5\,\Omega$ 产生的压降
\\
$5$ & $5\,\mathrm{V}$ 源 & 源的源值，方向决定它是加号还是减号
\tabbot
\end{tabular}

\subhdr{5.8 $I=1\,\mathrm{A}$（p37--p38）与功率自检}

把 $8\,\Omega$ 接回端口，戴维南等效给出

\[ I=\frac{U_{OC}}{R_0+R_L}=\frac{14}{6+8}=1\,\mathrm{A} , \]

与 p38 的 $I=1\,\mathrm{A}$ 一致。注意这个式子在课件里\textbf{从未写出}：p32 给了 $R_0$、p36 给了 $U_{OC}$，
p38 直接跳到 $I=1\,\mathrm{A}$（§10 记一条拼接缺口）。

\aside{代入本题的两个自检：（1）端口电压 $v=14-6\times1=8\,\mathrm{V}$，正好等于 $8\,\Omega\times1\,\mathrm{A}$；
（2）功率平衡：等效电源送出 $14\,\mathrm{V}\times1\,\mathrm{A}=14\,\mathrm{W}$，其中 $R_0$ 上耗 $1^2\times6=6\,\mathrm{W}$、
负载耗 $1^2\times8=8\,\mathrm{W}$，$6+8=14$ 对得上。}

再用原电路直接解一遍做交叉验证（把 $8\,\Omega$ 接上，两块仍只经公共线相连，各自的 KCL 多了一项负载电流）：

\[ V_a:\ \frac{V_a-8}{2}+\frac{V_a}{2}+1=0\Rightarrow V_a=3\,\mathrm{V};\qquad
   V_b:\ \frac{V_b-5}{5}+3-1=0\Rightarrow V_b=-5\,\mathrm{V} . \]

$V_a-V_b=8\,\mathrm{V}$，$I=8/8=1\,\mathrm{A}$，与等效法一致。表里对比空载与有载：

\par\vspace{2pt}
\noindent\begin{tabular}{@{}p{3.0cm}p{3.4cm}p{3.4cm}p{4.2cm}@{}}
\tabtop
\thead{状态} & \thead{$V_a$} & \thead{$V_b$} & \thead{端口电压 $V_a-V_b$}
\\
\tabmid
空载（p36） & $4\,\mathrm{V}$ & $-10\,\mathrm{V}$ & $14\,\mathrm{V}=U_{OC}$
\\
有载 $8\,\Omega$（p38） & $3\,\mathrm{V}$ & $-5\,\mathrm{V}$ & $8\,\mathrm{V}=8\,\Omega\times1\,\mathrm{A}$
\tabbot
\end{tabular}

%=============================================================
\subhdr{5.9 主例的两个变式问法（同一套等效参数，换着问）}

等效电路一旦求出，整族问题都能一行答完；反过来，凡是等效电路答不出的，必须回原电路。

\textbf{变式一：换负载。}$I=U_{OC}/(R_0+R_L)=14/(6+R_L)$。$R_L=6\,\Omega$ 时 $I=1.167\,\mathrm{A}$；
$R_L\to\infty$（开路）时 $v=U_{OC}=14\,\mathrm{V}$；$R_L=0$（短路）时 $I=2.33\,\mathrm{A}$（§6.5 直接验）。

\textbf{变式二：问内部量。}「$8\,\Omega$ 接上时 $3\,\mathrm{A}$ 电流源两端的电压是多少？」
等效电路里已经没有 $3\,\mathrm{A}$ 源，答不出；必须回原电路：有载时 $V_b=-5\,\mathrm{V}$（§5.8），
底母线为 $0\,\mathrm{V}$，故源端电压是 $-5\,\mathrm{V}$（上端低、下端高）。

判据很短：\textbf{问端口上的事，用等效电路；问内部的事，回原电路}。
混淆这两类问题，是「等效」这个词最常被误用之处。

%=============================================================
\hdr{§6}{p28 的「How to deduce this conclusion?」为什么推不出来}

\subhdr{6.1 课件做的只是「算了一次」}

p28 的版面是：左边写 $U_{OC}=U_S+R_1I_S$、$I_{SC}=U_S/R_1+I_S$，右边写
$R_0=U_{OC}/I_{SC}=(U_S+R_1I_S)/(U_S/R_1+I_S)=R_1$。这三步对\textbf{这一个}网络是对的，
但它不是证明，原因是最后一步「约掉」能成功纯属凑巧：

\[ U_S+R_1I_S=R_1\left(\frac{U_S}{R_1}+I_S\right), \]

分子恰好是分母的 $R_1$ 倍，所以比值必然等于 $R_1$。换句话说，p28 验证的是
「分子恰好是分母的常数倍」，而不是「$R_0=U_{OC}/I_{SC}$ 对任何线性二端网络都成立」。
换个网络（比如主例 p24 那个五支路网络），分子分母就没有这种倍数关系了，但等式照样成立——
这时候靠的是 §4.1 的直线，不是代数运气。

代入一组数看得更清楚：取 $U_S=12\,\mathrm{V}$、$R_1=4\,\Omega$、$I_S=3\,\mathrm{A}$，
则 $U_{OC}=12+12=24\,\mathrm{V}$、$I_{SC}=3+3=6\,\mathrm{A}$、$R_0=24/6=4\,\Omega=R_1$。
把 $I_S$ 改成 $1\,\mathrm{A}$：$U_{OC}=16\,\mathrm{V}$、$I_{SC}=4\,\mathrm{A}$，比值仍是 $4\,\Omega$。
无论怎么改，比值恒等于 $R_1$——因为分子本来就是分母的 $R_1$ 倍。所以 p28 验的是一个恒等式，不是一条定理。
真正的考验在主例上：那里 $U_{OC}=14\,\mathrm{V}$、$I_{SC}=7/3\,\mathrm{A}$，比值是 $6\,\Omega$，
而分母里并没有一个现成的「$6$」（§6.5）——这时只能靠 §4.1 的直线。

\subhdr{6.2 一般证明只有一步}

由 §4.1，端口方程是 $v=U_{OC}-R_0i_L$。短路就是令 $v=0$，此时

\[ i_L=\frac{U_{OC}}{R_0}\equiv I_{SC}\ \Longrightarrow\ R_0=\frac{U_{OC}}{I_{SC}} . \]

所以 $R_0=U_{OC}/I_{SC}$ 不是「又一条要背的公式」，而是直线方程的截距与斜率之比。
凡是「先分别求两个量、再相除」的题，背后都是这条直线；凡是直线不成立的地方，这个比值就没有意义（下一小节）。

\subhdr{6.3 什么时候 $U_{OC}/I_{SC}$ 失效}

\par\vspace{2pt}
\noindent\begin{tabular}{@{}p{4.0cm}p{4.4cm}p{5.6cm}@{}}
\tabtop
\thead{情形} & \thead{现象} & \thead{怎么办}
\\
\tabmid
$R_0=0$（端口等同理想电压源） & 短路电流趋于无穷，$U_{OC}/I_{SC}$ 变成「有限 $/\infty$」 & 直接说 $R_0=0$，等效电路就是一个电压源
\\
$R_0=\infty$（端口等同理想电流源） & 短路电流为 $0$，式子变成 $U_{OC}/0$ & 直接说 $R_0=\infty$，等效电路就是一个电流源
\\
$U_{OC}=0$ 且 $I_{SC}=0$ & 变成 $0/0$，无定义 & 端口是纯电阻（无源网络），$R_0$ 用置零法求
\\
网络含非线性元件 & 端口特性不是直线，两个数不足 & 叠加、戴维南同时失效，只能逐点求解
\tabbot
\end{tabular}

\subhdr{6.4 为什么能把网络切成两半（p34 的蓝框）}

p34 用浅蓝方框圈出「a、b 之间只经底部导线相通」。这句话是整道主例的合法性所在：
端口空载时，左块与右块唯一可能的电流通道被断开了（$8\,\Omega$ 已被拿掉），
公共线又是参考节点，所以\textbf{没有电流在左右两块之间流动}，两块各自满足「节点电流之和为零」，
可以各写一个 KCL。若两块之间还连着别的电阻或源，这个切法就不成立，
只能把整个网络当一个多节点问题解——这是「切两半」这一步的边界，课件没有点明。

%=============================================================
\subhdr{6.5 用主例直接验一次 $I_{SC}=U_{OC}/R_0$}

光有公式不算懂，用手算一次最短：把主例的 $8\,\Omega$ 换成一根导线（端口短路），此时
$V_a=V_b=V$（同一个节点）。左块的 KCL 多了一项短路电流 $I_{SC}$（从 a 端流出网络）：

\[ \frac{V-8}{2}+\frac{V}{2}+I_{SC}=0\ \Longrightarrow\ I_{SC}=4-V ; \]

右块的 KCL 里，短路电流从 b 端流入网络：

\[ \frac{V-5}{5}+3-I_{SC}=0 . \]

把 $I_{SC}=4-V$ 代进去：$\frac{V-5}{5}+3-4+V=0$，两边乘 $5$ 得 $V-5+5V-5=0$，
故 $V=\frac{5}{3}\,\mathrm{V}$、$I_{SC}=4-\frac{5}{3}=\frac{7}{3}\approx2.333\,\mathrm{A}$。
而 $\frac{U_{OC}}{R_0}=\frac{14}{6}=2.333\,\mathrm{A}$，两者一致——这条等式不是巧合，是 §4.1 那条直线在 $v=0$ 处的读数。

%=============================================================
\hdr{§7}{其余例题：一例半页，只写与主线不同的地方}

\subhdr{7.0 先选方法：四种题型对应四条路}

本册的两套工具（叠加、等效）加上第一章前面的节点法与支路电流法，一共四条路。做题时先认题型再动笔，
比「哪个公式熟用哪个」省时间。判据是从本册的例题里归纳出来的思路，不是新增考点：

\par\vspace{2pt}
\noindent\begin{tabular}{@{}p{5.4cm}p{3.6cm}p{5.4cm}@{}}
\tabtop
\thead{题目长什么样} & \thead{首选} & \thead{理由（本册出处）}
\\
\tabmid
只要一条支路的电流电压，其余部分源很多 & 戴维南，接回负载 & 一次端口等效，之后换负载都是一行（§5.9）
\\
所有支路的电流都要 & 节点法或支路电流法 & 叠加要解 $n$ 次，等效只顶替一条支路
\\
网络里带受控源 & 节点法；求 $R_0$ 可用置零法 & 受控源不能置零，但置零独立源后求 $R_0$ 仍成立（§2.4）
\\
电桥、桥式结构，串并联化不开 & 戴维南 + 比例法求 $R_0$ & 桥臂划出去就成串并联，短路电流最好算（§6.2）
\\
问功率 & 先求总电流，再平方 & 功率不是线性量，不能叠加（§3.3）
\tabbot
\end{tabular}

下面六个例子按这个判据各走一条路，答案与来路都写全，可以当自测。

\subhdr{7.1 E2（p16--p18）：$3\,\mathrm{A}$ 与 $5\,\mathrm{A}$ 为什么都取正}

p16 的电路是一个「田」字形网络。把元件写清：上支路 $7\,\mathrm{A}$ 电流源（箭头朝左，电流 T$\to$L）、
右上 $1\,\Omega$（T--R）、中 $2\,\Omega$（T--M）、中左 $3\,\Omega$（M--L）、中右 $35\,\mathrm{V}$ 源（$+$ 在 M 侧，M--R）、
下 $4\,\Omega$（L--R，$I$ 的参考方向 L$\to$R）。

\textbf{（1）只留 $7\,\mathrm{A}$（$35\,\mathrm{V}$ 短路）。}M 与 R 合并成一点，于是 $1\,\Omega\parallel2\,\Omega=2/3\,\Omega$（T 到 M/R），
$3\,\Omega\parallel4\,\Omega=12/7\,\Omega$（L 到 M/R）。$7\,\mathrm{A}$ 由 T 经源流到 L，再经 $12/7\,\Omega$ 回到 M/R，
经 $2/3\,\Omega$ 回到 T。L--R 这对节点上的电压是 $7\times12/7=12\,\mathrm{V}$，故

\[ I'=\frac{12\,\mathrm{V}}{4\,\Omega}=3\,\mathrm{A} , \]

方向由 L 到 R（也可直接用分流比：$7\times3/(3+4)=3\,\mathrm{A}$）。

\textbf{（2）只留 $35\,\mathrm{V}$（$7\,\mathrm{A}$ 开路）。}上支路在源处断开，T 只连着 $2\,\Omega$ 与 $1\,\Omega$，
二者串成 $3\,\Omega$ 并在 M--R 之间；另一条 M--L--R 是 $3+4=7\,\Omega$。两条都并在 $35\,\mathrm{V}$ 源两端，互不影响：

\[ I''=\frac{35\,\mathrm{V}}{7\,\Omega}=5\,\mathrm{A} , \]

路径是 M$\to$L$\to$R，与参考方向（L$\to$R）相同。

\textbf{（3）相加。}两个分量方向相同，都取正：$I=3+5=8\,\mathrm{A}$，与 p18 一致。
自检：有载时 $V_L-V_R=8\,\mathrm{A}\times4\,\Omega=32\,\mathrm{V}$。

顺便解释一件容易吓人的事：$35\,\mathrm{V}$ 源单独作用时，$2\,\Omega+1\,\Omega=3\,\Omega$ 那条路要吃
$35/3\approx11.7\,\mathrm{A}$，加上 $3\,\Omega+4\,\Omega=7\,\Omega$ 那条路的 $5\,\mathrm{A}$，
流过源的支路电流是 $16.7\,\mathrm{A}$（§7.4 的节点解给出同一个数）——远大于 $I''=5\,\mathrm{A}$。
这不是算错，是「源同时给两条并联路供电」。

\subhdr{7.2 E3（p19）：两个理想电流源都在时，节点方程只剩一条路}

课件这一页只给电路、不给答案。这里按文字层转录的拓扑解一遍，并把结论标 [待确认]（图面未获图像复核）。
节点取 L（左）、T（上）、M（中）、R（右）；取 $V_R=0$，底部 $10\,\mathrm{V}$ 源的 $+$ 在 L 侧，故 $V_L=10\,\mathrm{V}$。

T 上接着 $3\,\mathrm{A}$ 源（向 T）、$5\,\mathrm{A}$ 源（离 T）与 $4\,\Omega$ 一条路，所以 T 的 KCL 只经 $4\,\Omega$：

\[ \frac{V_T-V_M}{4}+5-3=0\ \Longrightarrow\ V_T-V_M=-8\,\mathrm{V} . \]

M 的 KCL（$2\,\Omega$ 到 L、$3\,\Omega$ 到 R、$4\,\Omega$ 到 T）：

\[ \frac{V_M-V_T}{4}+\frac{V_M-10}{2}+\frac{V_M}{3}=0\ \Longrightarrow\ V_M=3.6\,\mathrm{V},\quad V_T=-4.4\,\mathrm{V} . \]

再把 L 的 KCL 列出来：$3\,\mathrm{A}$ 经电流源离开 L，$(10-3.6)/2=3.2\,\mathrm{A}$ 经 $2\,\Omega$ 离开 L，
所以底部支路必须向 L 送入 $3+3.2=6.2\,\mathrm{A}$，即 $I=6.2\,\mathrm{A}$（参考方向自 R 向 L）。
这个例子要学的不是这个数，而是那句「T 的 KCL 只经 $4\,\Omega$」——两个理想电流源把 T 夹住了。

\subhdr{7.3 E4（p20）与 p40：两种算法给出同一个 $3\,\mathrm{A}$}

p20 与 p40 是同一个电路：$18\,\mathrm{V}$ 串 $1\,\Omega$、$9\,\mathrm{V}$ 串 $1\,\Omega$ 两条支路并接，向 $4\,\Omega$ 供电，
$I$ 的参考方向向下。

\textbf{（1）叠加。}$18\,\mathrm{V}$ 单独作用（$9\,\mathrm{V}$ 短路）：两只 $1\,\Omega$ 并联后与 $4\,\Omega$ 分压，
源电流 $18/(1+0.8)=10\,\mathrm{A}$，$4\,\Omega$ 分到 $10\times1/(1+4)=2\,\mathrm{A}$。
$9\,\mathrm{V}$ 单独作用同理得 $5\,\mathrm{A}$ 源电流、$4\,\Omega$ 分到 $1\,\mathrm{A}$。所以 $I=2+1=3\,\mathrm{A}$。

\textbf{（2）戴维南。}$U_{OC}=(18+9)/2=13.5\,\mathrm{V}$（两个 $1\,\Omega$ 把两个源平均分压），
$R_0=1\parallel1=0.5\,\Omega$，故 $I=13.5/(0.5+4)=3\,\mathrm{A}$。两条路一致。

\subhdr{7.4 用戴维南重做 E2（p39）：$R_0=3\,\Omega$、$U_{OC}=56\,\mathrm{V}$、$I=8\,\mathrm{A}$}

p39 的电路与 §7.1 的 E2 完全相同，只是要求改用戴维南。把 $4\,\Omega$ 划出去，端口取 a$=$L、b$=$R。

\textbf{（1）$R_0$：}$7\,\mathrm{A}$ 开路（上支路在源处断开）、$35\,\mathrm{V}$ 短路（M 与 R 合并）。
从 L 到 R 只剩 $3\,\Omega$；$2\,\Omega$ 串 $1\,\Omega$ 那一路两端都落在 R 上（M$\equiv$R、经 T 又回 R），
无源、无电流，不影响 $R_0$：

\[ R_0=3\,\Omega . \]

\aside{代入本题的判据：置零后 M 与 R 是同一个节点，$2\,\Omega+1\,\Omega$ 接成「从 R 出发又回到 R」的回路，
没有源就没有电流，对端口电阻没有贡献。这一步最容易被误算成 $3\parallel3=1.5\,\Omega$。}

\textbf{（2）$U_{OC}$：}取 $V_R=0$，则 $35\,\mathrm{V}$ 源把 $V_M$ 钉在 $35\,\mathrm{V}$。
T 的 KCL（$7\,\mathrm{A}$ 由 T 流向 L，故从 T 流出）：

\[ \frac{V_T-35}{2}+\frac{V_T-0}{1}+7=0\ \Longrightarrow\ V_T=7\,\mathrm{V} . \]

L 的 KCL（$7\,\mathrm{A}$ 流入 L，只从 $3\,\Omega$ 流出）：

\[ \frac{V_L-35}{3}=7\ \Longrightarrow\ V_L=56\,\mathrm{V},\qquad U_{OC}=56\,\mathrm{V} . \]

\textbf{（3）接回负载：}

\[ I=\frac{56}{3+4}=8\,\mathrm{A} , \]

与 §7.1 的叠加结果、与 p18 页的 $I=8\,\mathrm{A}$ 三者一致。这就是「$R_0$ 两种求法必须对得上」的实例。

\subhdr{7.5 一题四法（p42--p45）：四个答案必须一致}

p42--p45 是同一个电路贴四个方法名（「Source transformation」「Branch current analysis」
「Superposition theorem」「Thévenin's theorem」），\textbf{四页一个数都没有}。
电路：节点 A（上）、B（下）；左支路 $5\,\mathrm{A}$ 电流源（向上，注入 A）；中支路 $10\,\mathrm{V}$ 源（$+$ 在上）串 $2\,\Omega$；
右支路 $3\,\Omega$；$I_1$ 参考方向向上（中支路），$I_2$ 参考方向向下（$3\,\Omega$）。

\textbf{（1）电源变换。}$10\,\mathrm{V}$ 串 $2\,\Omega$ 变换成 $5\,\mathrm{A}\parallel2\,\Omega$（电流向上，与 $5\,\mathrm{A}$ 源同向），
两者合并成 $10\,\mathrm{A}$ 注入 A，并联电阻 $2\parallel3=1.2\,\Omega$：

\[ V_A=10\times1.2=12\,\mathrm{V},\qquad I_2=\frac{12}{3}=4\,\mathrm{A},\qquad
   I_1=\frac{10-12}{2}=-1\,\mathrm{A} . \]

\textbf{（2）支路电流法。}KCL：$5+I_1=I_2$；中支路 KVL：$V_A=10-2I_1$；右支路：$V_A=3I_2$。
三式联立得 $I_1=-1\,\mathrm{A}$、$I_2=4\,\mathrm{A}$、$V_A=12\,\mathrm{V}$。

\textbf{（3）叠加。}$5\,\mathrm{A}$ 单独作用（$10\,\mathrm{V}$ 短路）：$2\parallel3=1.2\,\Omega$，$V_A=6\,\mathrm{V}$，
$I_1'=-3\,\mathrm{A}$（中支路电流向下 $3\,\mathrm{A}$）、$I_2'=2\,\mathrm{A}$。
$10\,\mathrm{V}$ 单独作用（$5\,\mathrm{A}$ 开路）：回路 $2+3=5\,\Omega$、电流 $2\,\mathrm{A}$，源把电流顶进 A 再经 $3\,\Omega$ 下地，
$I_1''=+2\,\mathrm{A}$、$I_2''=+2\,\mathrm{A}$。相加：$I_1=-1\,\mathrm{A}$、$I_2=4\,\mathrm{A}$。

\textbf{（4）戴维南。}把 $3\,\Omega$ 划出去：$U_{OC}=10+5\times2=20\,\mathrm{V}$（$5\,\mathrm{A}$ 全部流过 $2\,\Omega$），
$R_0=2\,\Omega$（$5\,\mathrm{A}$ 开路、$10\,\mathrm{V}$ 短路）。接回：$V_A=20\times3/(2+3)=12\,\mathrm{V}$、$I_2=4\,\mathrm{A}$，
再由中支路 $I_1=(10-12)/2=-1\,\mathrm{A}$。

\par\vspace{2pt}
\noindent\begin{tabular}{@{}p{4.6cm}p{2.6cm}p{2.6cm}p{3.6cm}@{}}
\tabtop
\thead{方法} & \thead{$I_1$} & \thead{$I_2$} & \thead{中间量}
\\
\tabmid
电源变换 & $-1\,\mathrm{A}$ & $4\,\mathrm{A}$ & $V_A=12\,\mathrm{V}$
\\
支路电流法 & $-1\,\mathrm{A}$ & $4\,\mathrm{A}$ & $V_A=12\,\mathrm{V}$
\\
叠加定理 & $-1\,\mathrm{A}$ & $4\,\mathrm{A}$ & $V_A=6+6=12\,\mathrm{V}$
\\
戴维南定理 & $-1\,\mathrm{A}$ & $4\,\mathrm{A}$ & $U_{OC}=20\,\mathrm{V}$、$R_0=2\,\Omega$
\tabbot
\end{tabular}

四法结果相同。课件把这四页放在一起，本意大概就是「同一道题、四种工具、答案只有一个」；
可惜一页数字都没写，这一步的示范作用没发挥出来（§10 记一条）。

\subhdr{7.6 E5（p21）：一个「$3\,\Omega$ 上没有电流」的巧合}

同样只有电路、没有答案，拓扑按文字层转录复原并标 [待确认]。取左竖线为参考（$V_L=0$），
节点为 T、TR、M、BR；$3\,\mathrm{V}$ 源把 TR 与 BR 拉开 $3\,\mathrm{V}$。

T 的 KCL：$V_T/2+(V_T-V_{TR})/6-1=0$（$1\,\mathrm{A}$ 由 M 流向 T）；
M 的 KCL：$(V_M-V_{BR})/3+V_M/1+1=0$；TR 与 BR 的 KCL 相加消去 $3\,\mathrm{V}$ 源电流：
$(V_{TR}-V_T)/6+(V_{BR}-V_M)/3=0$；再加 $3\,\mathrm{V}$ 源的条件 $V_{TR}-V_{BR}=3$。

解得 $V_{BR}=-1\,\mathrm{V}$、$V_{TR}=2\,\mathrm{V}$、$V_T=2\,\mathrm{V}$、$V_M=-1\,\mathrm{V}$。
于是 $I$（参考方向 M$\to$L）为

\[ I=\frac{V_M-V_L}{1}=-1\,\mathrm{A} , \]

即实际有 $1\,\mathrm{A}$ 由 L 经 $1\,\Omega$ 流向 M。顺带一个巧合：$V_M=V_{BR}$ 使 $3\,\Omega$ 上电流为零，
$3\,\mathrm{V}$ 源支路也不出力——这类「某只电阻上恰好没电流」的结果是拓扑决定的，不是算错，
可以拿它当自检：把 $V_M$ 与 $V_{BR}$ 代回去，KCL 必须同时成立。

%=============================================================
\subhdr{7.7 p41 的习题 E3：为什么这一题留成练习}

p41 给的是 $12\,\mathrm{V}$ 串 $6\,\Omega$、一只并联 $6\,\Omega$、$1\,\Omega$ 与 $3\,\mathrm{A}$、以及待求支路 $5\,\Omega$（含 $I$）。
这一页的图在文字层里只剩元件名与相对位置，$1\,\Omega$ 与 $3\,\mathrm{A}$ 到底是「同一支路串联」还是「两路并接」，
两种读法给出的方程不同，而课件没有给答案可供反查（不像 E1、E2 那样有现成分量可以对拍）。
按本讲义的规矩，这种拓扑不确定的题标 [待确认]，不作为算例，只把方法固定下来：
把 $5\,\Omega$ 划出去当负载，求端口的 $U_{OC}$ 与 $R_0$，再接回 $\frac{U_{OC}}{R_0+5}$。
等你手上有清晰的图，按 §8.2 的四步走一遍即可；四步里没有一步依赖「图长什么样」之外的假设。

%=============================================================
\hdr{§8}{公式卡、四步清单与三处自检}

\subhdr{8.1 叠加：四步（按做题顺序）}

\begin{enumerate}
\item \textbf{定参考方向}：给待求量画箭头，全部分量沿用这个箭头（§3.2）；
\item \textbf{逐个置零}：只留一个独立源，其余电压源短路、电流源开路，受控源保留（§2.1、§2.2）；
\item \textbf{各算一次}：每个分电路用最省事的方法（分流、分压、单节点 KCL）求该分量，带符号记下；
\item \textbf{代数相加 + 自检}：$I=\sum I^{(k)}$；最后把结果代回原电路直接解一遍（§5.4）。
\end{enumerate}

\subhdr{8.2 戴维南/诺顿：四步}

\begin{enumerate}
\item \textbf{定端口}：把待求支路划出去，标 a、b；确认它是两块之间唯一的连接（§5.5）；
\item \textbf{求 $U_{OC}$}：端口空载，按拓扑解出 $V_a-V_b$（§5.7）；
\item \textbf{求 $R_0$}：内部独立源置零、受控源保留，从端口看进去算等效电阻；或用 $U_{OC}/I_{SC}$ 复核（§5.6、§6.2）；
\item \textbf{接回负载}：$I=U_{OC}/(R_0+R_L)$，端口电压 $v=U_{OC}-R_0I$；功率用总电流平方算（§3.3）。
\end{enumerate}

\par\vspace{2pt}
\noindent\begin{tabular}{@{}p{2.2cm}p{5.6cm}p{1.9cm}p{4.5cm}@{}}
\tabtop
\thead{量} & \thead{公式} & \thead{依据} & \thead{前提}
\\
\tabmid
$u^{(k)}$ & 只留第 $k$ 个源时的响应 & p2、p3 & 元件线性、受控源保留
\\
$u$ & $\sum_k u^{(k)}$ & p10 & 同一参考方向、代数相加
\\
$R_0$ & $U_{OC}/I_{SC}$ & p27 & 端口特性是直线；$R_0$ 有限且非零
\\
$I$ & $U_{OC}/(R_0+R_L)$ & §4.1 & 负载是纯电阻
\\
$v$ & $U_{OC}-R_0I$ & §4.1 & 同上
\\
$I_{SC}$ & $U_{OC}/R_0$ & p27 & 同上
\tabbot
\end{tabular}

\subhdr{8.3 最大功率传输 [教材补充]}

课件全册没有讲最大功率传输，但戴维南等效一出来，它就是一行求导：负载取多少时功率最大？
本讲义的依据是 §4.1 的等效电路，不是课件原文，故标 [教材补充]——补它的理由是：
「把网络化简成两个数」如果不落到「这两个数能回答什么问题」，学完就只剩背公式。

负载功率 $P=I^2R_L=\dfrac{U_{OC}^2R_L}{(R_0+R_L)^2}$。对 $R_L$ 求导并令其为零：

\[ \frac{\dd P}{\dd R_L}=U_{OC}^2\,\frac{R_0-R_L}{(R_0+R_L)^3}=0
   \ \Longrightarrow\ R_L=R_0,\qquad P_{\max}=\frac{U_{OC}^2}{4R_0} . \]

代主例的数：$U_{OC}=14\,\mathrm{V}$、$R_0=6\,\Omega$，故 $R_L=6\,\Omega$ 时
$P_{\max}=196/24=8.17\,\mathrm{W}$；课件 p38 取的 $8\,\Omega$ 负载给出 $1^2\times8=8.00\,\mathrm{W}$，
只差 $2\%$——p38 那个负载其实是「接近最优」的取值。

\aside{代入本题的核对：$R_L=6\,\Omega$ 时 $I=14/12=1.167\,\mathrm{A}$，$P=1.167^2\times6=8.17\,\mathrm{W}$；
$R_L=8\,\Omega$ 时 $I=1\,\mathrm{A}$，$P=8\,\mathrm{W}$；$R_L=4\,\Omega$ 时 $I=1.4\,\mathrm{A}$，$P=7.84\,\mathrm{W}$。
三个数说明功率对 $R_L$ 在 $R_0$ 附近很平缓，这正是「匹配」在实际中不苛求精确的原因。}

\par\vspace{2pt}
\noindent\begin{tabular}{@{}p{2.4cm}p{2.6cm}p{2.6cm}p{2.6cm}p{3.6cm}@{}}
\tabtop
\thead{$R_L$} & \thead{$I$} & \thead{$v=U_{OC}-R_0I$} & \thead{$P=I^2R_L$} & \thead{备注}
\\
\tabmid
$0$（短路） & $2.33\,\mathrm{A}$ & $0$ & $0$ & $I=I_{SC}$，功率全落在 $R_0$ 上
\\
$3\,\Omega$ & $1.56\,\mathrm{A}$ & $4.67\,\mathrm{V}$ & $7.26\,\mathrm{W}$ & 偏离匹配点
\\
$6\,\Omega$ & $1.17\,\mathrm{A}$ & $7.00\,\mathrm{V}$ & $8.17\,\mathrm{W}$ & 匹配，$P$ 最大
\\
$8\,\Omega$ & $1.00\,\mathrm{A}$ & $8.00\,\mathrm{V}$ & $8.00\,\mathrm{W}$ & 课件 p38 的取值
\\
$12\,\Omega$ & $0.78\,\mathrm{A}$ & $9.33\,\mathrm{V}$ & $7.26\,\mathrm{W}$ & 与 $3\,\Omega$ 对称
\\
$\infty$（开路） & $0$ & $14\,\mathrm{V}$ & $0$ & $v=U_{OC}$
\tabbot
\end{tabular}

这张表就是「两个数能回答哪些问题」的清单：给定 $R_L$，$I$、$v$、$P$ 全部一行算出。

\subhdr{8.4 三处自检（每题都该做）}

\begin{enumerate}
\item \textbf{分量回代}：把结果代回原电路直接解一遍，与叠加/等效的结果比（§5.4、§5.8）；
\item \textbf{两个 $R_0$ 对拍}：置零法与 $U_{OC}/I_{SC}$ 必须相等（§7.4）；
\item \textbf{功率平衡}：等效电源送出的功率 $=$ 内阻损耗 $+$ 负载功率（§5.8：$14=6+8$）。
\end{enumerate}

%=============================================================
\hdr{§9}{易错点（每条都能指回某页某一步）}

\gapitem{（1）把功率当线性量}{p11 第 (2) 条只给结论。E1 的 $10\,\Omega$ 上：$P'+P''=40+10=50\,\mathrm{W}$，
真值只有 $10\,\mathrm{W}$，差的是交叉项 $-40\,\mathrm{W}$（§3.3）。凡问功率，先求总电流。}

\gapitem{（2）分量不带符号}{E1 的 $I''=-1\,\mathrm{A}$（§5.3）。若提前把箭头翻过来，会得 $3\,\mathrm{A}$；
分量算出负数时必须保留负号，负号说明实际方向与参考方向相反。}

\gapitem{（3）把受控源也置零}{受控源是线性元件的一部分，置零它等于把放大能力删掉（§2.2）。
课件 Notice 第 (1) 条只提「元件线性」，没提这一条，但电子电路里一旦置零受控源，结果必错。}

\gapitem{（4）求 $R_0$ 时把负载算进去}{p31 的 $R_0$ 是「$8\,\Omega$ 已划出去之后」的端口电阻；
若把 $8\,\Omega$ 并进去，得到的是「端口总电阻」而不是定理里的 $R_0$（§5.6、§0.3）。}

\gapitem{（5）求 $U_{OC}$ 时端口没断开}{$U_{OC}$ 的定义是端口开路（$I=0$）。带着负载量到的电压是 $8\,\mathrm{V}$（§5.8 表），
不是 $14\,\mathrm{V}$；两者差 $6\,\Omega\times1\,\mathrm{A}$。}

\gapitem{（6）$R_0$ 的串并联搞反}{主例里 $R_0=(2\parallel2)+5=6\,\Omega$，不是 $(2\parallel2)\parallel5$。
判据是「a 到 b 的通路必须串穿两块」（§5.6）。}

\gapitem{（7）以为「等效」是整体替换}{等效只对\textbf{端口外}成立：换掉的那部分网络内部各元件的电流电压全都变了
（p38 的等效电路里已经没有 $8\,\mathrm{V}$、$3\,\mathrm{A}$、$5\,\mathrm{V}$ 这些源）。题目若还问内部量，不能先等效再回头算。}

\gapitem{（8）$U_{OC}/I_{SC}$ 用在 $R_0=0$ 或 $R_0=\infty$ 上}{此时比值无意义（§6.3）。
先判断端口是「理想电压源型」还是「理想电流源型」，再决定用哪种求法。}

\gapitem{（9）在非线性元件上用叠加}{二极管、稳压管、白炽灯丝都不满足「电导与源值无关」（§1.2）。
课件 p2 的 「linear electric circuit」 是使用前提，不是修饰语。}

\gapitem{（10）把齐次性当万能结论}{「源值翻倍、响应翻倍」只对线性量成立（§1.5）。
功率、二极管的电流、三极管的 $\beta$ 都不是这样；翻倍之前先问一句「这个量对源值是一次函数吗」。}

\gapitem{（11）照抄 p11 第 (3) 条}{第 (3) 条标题说方向、正文讲功率（§3.1）。照抄下来，笔记里就整块缺掉「方向」，
而它恰恰是 E1 里 $I''=-1\,\mathrm{A}$ 那一步的全部理由。抄笔记时按标题补正文，不按正文抄标题。}

%=============================================================
\hdr{§10}{缺口清单：课件哪句话跳了步，本讲义补在哪}

\gapitem{机制缺口}{p2、p3\quad 「the algebraic sum of the currents … when each source acts alone」\quad
缺：为什么「分别作用再相加」等于「一起作用」（线性性）\quad 补：§1.2 用 $\mathbf{G}\mathbf{v}=\mathbf{J}$ 两行证明}

\gapitem{拼接缺口}{p3 与 p5\quad 「replaced by their internal resistances」对「short circuit」/「open circuit」\quad
缺：两种说法之间的桥（置零置的是源值，内阻留下）\quad 补：§2.1 对照表}

\gapitem{机制缺口}{p5--p9\quad 「Assume only this one is working」\quad
缺：为什么可以「只让一个源工作」（合法性来自线性，不是电路操作）\quad 补：§2.2}

\gapitem{适用前提}{p5--p9、p11\quad 全册未提受控源\quad
缺：受控源不能置零、必须留在每个分电路里\quad 补：§2.2 的 warnbox、§9 第 (3) 条}

\gapitem{形式改写}{p11 第 (3) 条\quad 「Pay attention to the direction」 的正文整段重复第 (2) 条\quad
缺：方向规则本身（分量带符号、不翻箭头）\quad 补：§3.2}

\gapitem{代数跳步}{p11 第 (2) 条\quad 「cannot be determined by a simple sum of the powers」\quad
缺：差多少、差在哪（交叉项 $2I'I''R$）\quad 补：§3.3 用 E1 的数字算 $40+10-40=10$}

\gapitem{代数跳步}{p14--p15\quad 「$I'=2\,\mathrm{A}$」「$I''=-1\,\mathrm{A}$」\quad
缺：两个分量怎么算、$-1$ 的负号从哪来\quad 补：§5.2 分流比、§5.3 方向分析}

\gapitem{代数跳步}{p18\quad 「$I'=3\,\mathrm{A}$」「$I''=5\,\mathrm{A}$」\quad
缺：$3\,\mathrm{A}$ 与 $5\,\mathrm{A}$ 的算法，以及为什么两个都取正\quad 补：§7.1}

\gapitem{机制缺口}{p22--p23\quad 「can be replaced by an equivalent circuit」\quad
缺：为什么外部行为只要两个数（端口特性是直线），以及定理的证明\quad 补：§4.1（用叠加证明，标 [教材补充]）}

\gapitem{形式改写}{p26、p27\quad 「$U_{eS}=U_{OC}$」、$R_0=U_{eS}/I_{SC}$\quad
缺：两条都是原定理的推论，而不是定理本身\quad 补：§4.2 由 $v=U_{OC}-R_0i_L$ 推出}

\gapitem{代数跳步}{p28\quad 「How to deduce this conclusion?」\quad
缺：课件只对本网络算了一次（分子恰是分母的 $R_1$ 倍），没有一般证明\quad 补：§6.1、§6.2}

\gapitem{代数跳步}{p36\quad 「$U_{ab}+5-15-4=0$」\quad
缺：$4$、$15$、$5$ 三个数各自从哪来\quad 补：§5.7 第 (4) 点的来源表}

\gapitem{拼接缺口}{p32、p36 与 p38\quad $R_0=6\,\Omega$、$U_{ab}=14\,\mathrm{V}$ 之后直接 「$I=1\,\mathrm{A}$」\quad
缺：中间那一步 $I=U_{OC}/(R_0+R_L)$ 从未写出\quad 补：§5.8}

\gapitem{答案形态}{p19、p21、p39--p45\quad 页面上只有「E3」「E5」「Method 1-4」这几个标签，没有任何结果\quad
缺：E3、E5 与四法四页全无结果\quad 补：§7.2、§7.6、§7.5（E3、E5 的拓扑标 [待确认]）}

\gapitem{机制缺口}{p1--p21\quad 每页页眉都是「Superposition theorem」，正文从不提「齐次」\quad
缺：线性 $=$ 齐次 $+$ 可加，齐次性是验算响应的捷径\quad 补：§1.5 用 E1 的 $1\,\mathrm{A}\to2\,\mathrm{A}$ 验一次}

\gapitem{适用区间}{p22\quad 「active linear two-terminal network」\quad
缺：非线性网络为什么不能只用两个数（折线不是直线）\quad 补：§4.5}

\gapitem{版面}{p5--p9 每页都重复同一行「Content:」与同一张图（五帧）、p29--p38 同一电路十帧、p42--p45 同一电路四页\quad
缺：内容只应计一次\quad 补：§0 坐标表按去重计，此处备案}

%=============================================================
\hdr{§11}{三句话总结}

\begin{enumerate}
\item \textbf{叠加是解的结构，不是操作}：线性元件让 $\mathbf{G}$ 与源值无关，于是 $\mathbf{v}=\sum_k\mathbf{v}_k$；
      分量必须带符号，功率必须用总电流重算。
\item \textbf{戴维南是叠加的一次应用}：端口特性是直线 $v=U_{OC}-R_0i_L$，两个数就够；
      $R_0=U_{OC}/I_{SC}$ 是这条直线的截距比，不是新公式。
\item \textbf{两个主例题各给一条纪律}：E1 教「分量算完必须带符号相加」，
      主例教「先划端口、再置零求 $R_0$、最后接回负载」，每一步都能用原电路直接解来对拍。
\end{enumerate}

%=============================================================
\hdr{§12}{与课件的页码对照表}

\par\vspace{2pt}
\noindent\begin{tabular}{@{}p{1.7cm}p{6.6cm}p{6.4cm}@{}}
\tabtop
\thead{课件页} & \thead{课件内容} & \thead{本讲义位置}
\\
\tabmid
p1--p2 & 1.9 封面、定义与用途 & §1.1
\\
p3--p4 & 正式表述、通用网络图、两处提问 & §1.1、§2.1
\\
p5--p9 & 置零规则与五帧动画 & §2.1、§2.2
\\
p10 & 三步法 & §2.3、§8.1
\\
p11 & Notice 四条（第 (3) 条为空） & §3
\\
p12--p15 & E1：$2\,\mathrm{A}$、$-1\,\mathrm{A}$、$1\,\mathrm{A}$ & §5.1--§5.4
\\
p16--p18 & E2：$3\,\mathrm{A}$、$5\,\mathrm{A}$、$8\,\mathrm{A}$ & §7.1、§7.4
\\
p19 & E3（无答案） & §7.2
\\
p20 & E4（无答案） & §7.3
\\
p21 & E5（无答案） & §7.6
\\
p22 & 1.10.1 术语、分类、两幅等效图 & §4.0、§4.3
\\
p23 & 戴维南定理陈述 & §4.1
\\
p24--p25 & 主例原电路、划端口 & §5.5
\\
p26--p27 & $U_{eS}=U_{OC}$、$R_0=U_{eS}/I_{SC}$ & §4.2、§4.4
\\
p28 & 「How to deduce this conclusion?」 & §6.1、§6.2
\\
p29--p32 & Terminal $\to$ Equivalent resistance $\to 6\,\Omega$ & §5.6
\\
p33--p36 & Open-circuit voltage $\to 14\,\mathrm{V}$ & §5.7
\\
p37--p38 & 等效电路 $\to 1\,\mathrm{A}$ & §5.8、§8.3
\\
p39--p41 & 习题 E1--E3 & §7.4、§7.3、§10
\\
p42--p45 & E4 与四法标签 & §7.5
\tabbot
\end{tabular}

\flowerrule

{\small\color{inkgray}
本讲义所有数值已用 Python 独立复算（E1：$I'=2\,\mathrm{A}$、$I''=-1\,\mathrm{A}$、$I=1\,\mathrm{A}$；
E2：$3\,\mathrm{A}$、$5\,\mathrm{A}$、$8\,\mathrm{A}$；主例：$R_0=6\,\Omega$、$U_{OC}=14\,\mathrm{V}$、$I=1\,\mathrm{A}$，有载 $V_a=3\,\mathrm{V}$、$V_b=-5\,\mathrm{V}$；
E4：$3\,\mathrm{A}$；p39：$R_0=3\,\Omega$、$U_{OC}=56\,\mathrm{V}$、$I=8\,\mathrm{A}$；
四法：$I_1=-1\,\mathrm{A}$、$I_2=4\,\mathrm{A}$、$V_A=12\,\mathrm{V}$；最大功率：$8.17\,\mathrm{W}$）。
标 [教材补充] 的两处（戴维南定理的叠加证明、最大功率传输）不是课件原文，只作理解用，不作文档口径。
E3（§7.2）与 E5（§7.6）的电路拓扑按文字层转录复原，图面未获图像复核，结论标 [待确认]。}

\end{document}
